\subsection{Fuentes oficiales}

El sistema integra sensores remotos, geometría, clima, topografía y estadísticas públicas.
La \cref{tab:cp01-sources} resume el inventario principal congelado en el contrato de
datos.

\begin{longtable}{p{0.20\linewidth}p{0.19\linewidth}p{0.49\linewidth}}
\caption{Fuentes principales y semántica operacional.}
\label{tab:cp01-sources}\\
\toprule
Fuente & Cobertura & Uso principal \\
\midrule
\endfirsthead
\toprule
Fuente & Cobertura & Uso principal \\
\midrule
\endhead
Split oficial & 197 parcelas & ID, área, conjunto y rendimiento observado para 138 parcelas.\\
Geometrías oficiales & 197 polígonos, EPSG:4326 & centroides, distancias, extracción espacial y QA administrativa.\\
BASIC & $107{,}666\times96$; 2022--2025 & Sentinel-2/Landsat, fenología histórica y 2025, índices de vegetación/agua/estrés.\\
PRO & $47{,}804\times20$; 2025 & Planet NDVI/EVI/LAI/MSAVI y estadísticos temporales.\\
Clima oficial & 48 meses por variable, 2022--2025 & precipitación, Tmin, Tmax, climatología/anomalías.\\
Topografía oficial & 120 m & elevación y pendiente.\\
INEGI municipios & Hidalgo/Puebla/Tlaxcala & CVEGEO, unión SIAP y QA espacial.\\
CHIRPS v3 & 214 días, abr--oct 2025 & precipitación diaria y episodios secos/extremos; excluido en varias fases por QC.\\
WaPOR v3 & 21 dekads, abr--oct 2025 & AETI y NPP; balance hídrico/productividad.\\
SoilGrids & 9 propiedades $\times$ 4 profundidades & perfil edáfico estático y heterogeneidad vertical.\\
SIAP & 2003--2025 & contexto municipal histórico y proxy contemporáneo 2025.\\
INEGI CEM & $\sim15$ m & elevación/relieve de mayor resolución.\\
\bottomrule
\end{longtable}

\subsection{BASIC: representación longitudinal multiespectral}

BASIC contiene 107,666 registros y 96 columnas. Los registros cubren las 197 parcelas entre
2022-01-01 y 2025-12-31 y provienen de Landsat y Sentinel-2. Las fechas se interpretan
estrictamente como \artifact{DD/MM/YYYY}; la llave parcela--fecha--sensor es única.

La estructura de variables comprende 23 familias de índices con cuatro resúmenes espaciales
por observación ---promedio, desviación estándar, máximo y mínimo---, además de metadatos de
fecha, sensor y nubosidad. Dos decisiones son esenciales:
\begin{enumerate}
\item no se mezclan silenciosamente variables homónimas entre sensores;
\item la ausencia sistemática de ciertos índices por sensor no se corrige con un fill global,
porque esa ausencia es estructural y contiene provenance.
\end{enumerate}

En particular, fórmulas provider-specific como MRA y VI6T no se inventan. Se conserva su
semántica tal como aparece en el suministro original. Esta política es importante porque un
índice de nombre opaco no debe recibir una interpretación espectral creada a posteriori sólo
para justificar una feature.

\subsection{PRO: señal Planet 2025}

PRO aporta 47,804 registros por 20 columnas, todas las parcelas y fechas de 2025-01-01 a
2025-12-30. El sensor Planet contiene NDVI, EVI, LAI y MSAVI con estadísticos espaciales.
Esta fuente es particularmente relevante para capturar el ciclo contemporáneo con mayor
densidad/resolución, pero sigue siendo una colección de observaciones longitudinales. Su
incorporación supervisada ocurre sólo después de agregarla de forma parcelaria o de construir
métricas temporales que respetan la entidad parcela.

La coexistencia de Sentinel-2 y Planet permite además construir features de concordancia
inter-sensor. Esa comparación no presupone equivalencia radiométrica: se utiliza para
cuantificar si ambas cadenas describen una forma temporal compatible para la misma parcela.

\subsection{Clima oficial, CHIRPS y WaPOR}

La fuente climática oficial contiene 48 rasters mensuales para cada una de precipitación,
temperatura mínima y temperatura máxima, con continuidad 2022--2025 y CRS EPSG:4326. A partir
de ella se construyen climatologías históricas, totales estacionales y anomalías 2025.

CHIRPS v3 aporta precipitación diaria para 214 fechas entre abril y octubre de 2025. Aunque su
resolución temporal permite construir número de días de lluvia, máximos móviles y duración de
episodios secos, el pipeline mantuvo CHIRPS fuera de varias representaciones de modelado porque
Feature Table v1 sólo retuvo una característica y se decidió no sobreinterpretar una
extracción que requería QC adicional.

WaPOR v3 aporta AETI y NPP en 21 periodos dekadales. Dado que sus bandas representan tasas,
una agregación estacional correcta multiplica por la duración del dekad:
\begin{equation}
AETI_{\mathrm{season}}
=
\sum_{k=1}^{21} AETI_k d_k.
\label{eq:wapor-duration}
\end{equation}
La misma regla se aplica a otros productos de tasa. Esta distinción impide sumar valores con
unidades por día como si fueran acumulados. La capa agronómica utilizará posteriormente estas
magnitudes para construir proxies de productividad hídrica \citep{Dhonthi2024}.

\subsection{Suelo y topografía}

SoilGrids aporta 36 GeoTIFF: nueve propiedades por cuatro profundidades:
pH, arcilla, arena, limo, carbono orgánico, nitrógeno, capacidad de intercambio catiónico,
densidad aparente y fracción de fragmentos gruesos. Los TIFF descargados mediante WCS pueden
carecer de CRS en el header; sólo los productos positivamente identificados reciben el
fallback ESRI:54052 definido por el contrato. Las conversiones de escala INT16 a unidades
físicas se aplican según la especificación del producto, nunca por inferencia ad hoc.

La topografía combina el producto oficial de 120 m y el CEM estatal de INEGI de
aproximadamente 15 m. Las fuentes se mantienen separadas mediante prefijos para evitar que una
resolución sustituya silenciosamente a otra. En modelos locales, la elevación y la pendiente
pueden actuar como contexto ambiental; sin embargo, la distancia geográfica directa terminó
teniendo una asociación más fuerte con diferencias de rendimiento que la distancia
multivariada completa.

\subsection{SIAP y la llave administrativa}

El archivo histórico SIAP cubre 2003--2025. La llave municipal estable es
\begin{equation}
\mathrm{CVEGEO}
=
\mathrm{zfill}_2(\mathrm{Idestado})
\Vert
\mathrm{zfill}_3(\mathrm{Idmunicipio}),
\label{eq:cvegeo}
\end{equation}
donde $\Vert$ denota concatenación. Nunca se une únicamente por
\artifact{Idmunicipio}, porque dicho identificador no es globalmente único entre estados.

Se identificaron al menos las categorías \artifact{Cebada grano} y
\artifact{Cebada forrajera en verde}, además de ciclos Otoño--Invierno y
Primavera--Verano y modalidades Riego/Temporal. Mezclar grano y forraje modificaría la
semántica de la respuesta, por lo que se prohíbe explícitamente.

En Feature Table v1 se preservó una versión histórica del agregador SIAP. En
Checkpoint~04E.1 se volvió a la tabla detallada para filtrar de forma explícita
\artifact{Cebada grano + Primavera-Verano + Temporal + CVEGEO + 2025}.
La distinción entre el agregado legado y el scope exacto es importante para interpretar por
qué una fuente administrativamente correcta puede seguir siendo un proxy pobre del rendimiento
parcelario. Debe señalarse además que la elección de Primavera--Verano se alineó con el ciclo
abril--octubre del reto; no se presenta aquí como una certificación externa independiente de
semántica de ciclo.

\subsection{Asignación municipal y casos de borde}

La etiqueta administrativa oficial de parcela se preserva como referencia para joins. La
superposición con polígonos INEGI se utiliza como QA/fallback y no sustituye silenciosamente
el dato oficial. Se documentaron, entre otros, AGC\_020 con una pequeña intersección
interestatal, AGC\_129 atravesando municipios en Hidalgo y AGC\_048 con discrepancia entre
etiqueta oficial y overlay. Esta política evita que un algoritmo de overlay convierta
automáticamente un caso geométrico limítrofe en una reasignación administrativa.

\subsection{Fuentes investigadas pero no promovidas}

ERA5-Land fue explorado y se implementó tooling de descarga, pero las solicitudes a CDS
presentaron problemas repetidos de costo/backend y no se obtuvo un artefacto completo
certificado para el pipeline. En lugar de permitir una fuente parcial, se declaró opcional y
se excluyó de downstream. Esta decisión ilustra una regla general: una fuente potencialmente
útil no entra al modelo por prestigio o disponibilidad nominal; debe tener cobertura,
unidades y proceso de extracción auditables.
